% !TeX root = ../supplementary.tex
\section{Camera Annotation Details}
\label{sec:supp_camera}

The main paper presents the comparisons that determine our camera-annotation pipeline.
Here we provide the component ablation, the full scale-recovery comparison, and implementation details.

\subsection{Hybrid reconstruction}

Table~\ref{tab:supp_hybrid} evaluates native Pi3X~\cite{wang2026pi3permutationequivariantvisualgeometry}, shared-intrinsics BA, and a direct UniDepthV2~\cite{piccinelli2025unidepthv2} scale substitution on an eight-clip subset.
BA reduces median Sim3 ATE from 0.2079 to 0.1358 and scale error from 38.87\% to 32.95\%, with a small increase in rotation error.
Direct depth substitution before the later calibration leaves 36.48\% scale error.
These results motivate the combination of learned initialization, geometric refinement, and a calibrated scale stage.

\begin{table}[htbp]
    \centering
    \caption{Components of the hybrid camera pipeline. All configurations are valid on the same eight clips. Entries are medians, with Sim3 ATE in reference-trajectory units.}
    \label{tab:supp_hybrid}
    \small
    \setlength{\tabcolsep}{3pt}
    \begin{tabular}{lrrr}
        \toprule
        Configuration & Sim3 ATE $\downarrow$ & Rot. ($^\circ$) $\downarrow$ & Scale (\%) $\downarrow$ \\
        \midrule
        Pi3X, native & 0.2079 & \textbf{0.4252} & 38.87 \\
        + shared-$K$ BA & \textbf{0.1358} & 0.4422 & \textbf{32.95} \\
        \quad + UniDepthV2 scale & \textbf{0.1358} & 0.4422 & 36.48 \\
        \bottomrule
    \end{tabular}
\end{table}

\subsection{Learned geometry comparison}

Table~\ref{tab:supp_learned_geometry} compares relative geometry and recovered scale.
ViPE~\cite{huang2025vipevideoposeengine} uses its \texttt{pose\_only} profile.
Both Pi3X and ViPE use a common UniDepthV2 scale backend.
The native-ViPE branch retains ViPE's original geometry and was evaluated as a post-hoc diagnostic.
Coverage uses all 16 attempted clips, while error statistics use the 14 common successful outputs.

\begin{table}[htbp]
    \centering
    \caption{Relative geometry and scale. Error medians use 14 common successful clips. Sim3 ATE is normalized by reference path length, while coverage uses all 16 clips.}
    \label{tab:supp_learned_geometry}
    \small
    \setlength{\tabcolsep}{3pt}
    \begin{tabular}{lrrrr}
        \toprule
        Geometry & Full & ATE (\%) $\downarrow$ & Rot. ($^\circ$) $\downarrow$ & Scale (\%) $\downarrow$ \\
        \midrule
        Pi3X + BA & \textbf{16/16} & 0.127 & 0.460 & 22.17 \\
        ViPE + BA & 14/16 & 0.125 & 0.459 & \textbf{20.64} \\
        ViPE, native & 14/16 & \textbf{0.083} & \textbf{0.455} & 21.30 \\
        \bottomrule
    \end{tabular}
\end{table}

\subsection{Scale-recovery selection}

Table~\ref{tab:supp_scale} includes crop augmentation, MoGe-2~\cite{wang2025moge2accuratemonoculargeometry}, and depth-model routing in addition to the baseline and selected inverse-depth correction.
The correction achieves the lowest median scale error, while MoGe-2 and routing have lower P90 errors.
We retain MoGe-2 as a recovery option.
The residual median and tail errors motivate combining inferred annotations with directly recorded geometry and self-captured UAV observations.

\begin{table}[htbp]
    \centering
    \caption{Full scale-recovery comparison. Geometry and intrinsics are fixed. Each frozen configuration has 15 scorable outputs from 16 attempted clips.}
    \label{tab:supp_scale}
    \small
    \setlength{\tabcolsep}{3pt}
    \begin{tabular}{lrr}
        \toprule
        Scale configuration & Median (\%) $\downarrow$ & P90 (\%) $\downarrow$ \\
        \midrule
        UniDepthV2, P98 baseline & 25.46 & 51.65 \\
        + inverse-depth correction & \textbf{22.40} & 45.60 \\
        Crop augmentation & 26.77 & 53.15 \\
        MoGe-2, P98 & 24.04 & 44.65 \\
        UniDepthV2 / MoGe-2 routing & 25.05 & \textbf{44.22} \\
        \bottomrule
    \end{tabular}
\end{table}

\begin{table*}[t]
    \centering
    \caption{Depth annotation with focal-input and temporal-stabilization variants on 16 GTA clips. Metrics are averaged equally across clips. Lower is better.}
    \label{tab:supp_depth}
    \small
    \setlength{\tabcolsep}{8pt}
    \begin{tabular}{@{}llrrrr@{}}
        \toprule
        & & \multicolumn{2}{c}{AbsRel (\%)} & RMSE (m) & Temporal \\
        \cmidrule(lr){3-4}
        Configuration & External focal & 0.1--80\,m & 0.1--1,000\,m & 0.1--80\,m & log error \\
        \midrule
        MetricAnything & Yes & 20.61 & 23.80 & 10.57 & 0.03778 \\
        MetricAnything & No & \textbf{17.69} & \textbf{21.47} & \textbf{9.54} & 0.03759 \\
        MetricAnything + VDA & Yes & 20.93 & 24.03 & 10.72 & 0.02210 \\
        MetricAnything + VDA & No & 18.20 & 21.86 & 9.63 & \textbf{0.02124} \\
        \bottomrule
    \end{tabular}
\end{table*}

\subsection{Geometric refinement and scale estimation}

We transform initialized poses to the first-frame coordinate system and use a shared pinhole camera with a centered principal point.
The initial shared focal length is the median of the predicted horizontal and vertical focal lengths across frames.
Multiview feature tracks support triangulation and BA of poses and scene points.
The shared focal length is refined when geometric observability and held-out reprojection checks permit it.
A rejected focal update triggers BA with the initial focal length fixed.

Scale estimation uses UniDepthV2 predictions on 32 frames distributed across the clip.
Reliable reconstructed points provide depth anchors, balanced over a $4\!\times\!3$ spatial grid and weighted by geometric reliability and predicted depth uncertainty.
For each frame, we robustly fit log-depth ratios between the model predictions and reconstructed optical-axis depths.
We reject inconsistent frames, apply median-absolute-deviation filtering to frame scales, and require a retained P90/P10 ratio of at most 2.
The 98th percentile of the retained frame scales gives the P98 baseline.

The selected correction fits an affine relation between normalized inverse depths, with slope constrained to $[0.5,2]$ and offset to $[-0.3,0.3]$.
It is accepted when held-out points improve by at least 5\% in at least four independently validated frames.
Only the fitted slope contributes to scale conversion.
If these checks fail, the pipeline retains the baseline scale.
One positive scalar rescales all camera translations and reconstructed points in the clip.

\section{Depth Annotation Details}
\label{sec:supp_depth}

\subsection{Models and temporal scale stabilization}

MetricAnything Student PointMap~\cite{ma2026metricanything} predicts per-frame optical-axis depth using resolution level 9 and its internally recovered focal length.
The metric Large (ViT-L) Video Depth Anything model~\cite{chen2025videodepthanything} runs at input size 518 with 32-frame temporal windows and 10-frame overlap.
Both receive RGB frames in the selected configuration.

Let $M_t(p)$ and $V_t(p)$ denote MetricAnything and VDA depth at pixel $p$ in frame $t$.
We compute a robust log-depth ratio on a valid anchor set $\Omega_t$,
\begin{equation}
    r_t = \operatorname{median}_{p\in\Omega_t}
    \bigl(\log M_t(p)-\log V_t(p)\bigr),
\end{equation}
then stabilize MetricAnything with the clip-wide median ratio,
\begin{equation}
    \bar r = \operatorname{median}_t r_t,
    \qquad D_t(p) = M_t(p)\exp(\bar r-r_t).
\end{equation}
The framewise scalar preserves spatial structure while following VDA's relative temporal scale.
Anchors are sampled every eight pixels within the central region of each $832\!\times\!480$ frame, excluding 40 horizontal and 24 vertical border pixels.
We retain finite, non-sky predictions for which both depths lie between 1 and 150\,m.
Every frame must contain an anchor and the clip must contain more than 100 anchors in total.
Clips with insufficient support retain the original MetricAnything predictions.

\subsection{Sky processing and depth encoding}

The sky head of DA3METRIC-LARGE~\cite{lin2025depthanything3} operates at resolution 504.
We align neighboring sky scores with bidirectional DIS optical flow and take a three-frame temporal median, keeping the endpoint frames unchanged.
Correspondences with forward--backward error above 1.5 pixels at the target resolution are rejected.
The stabilized sky-score threshold is 0.7.

We downsample metric depth to $448\!\times\!256$ using validity-weighted area averaging.
A target pixel is accepted as non-sky when its non-sky area coverage is at least $1-10^{-6}$.
Depth is clipped to $[0.1,2000]$\,m and mapped to the shared logarithmic range,
\begin{equation}
    \ell(z) = 2\frac{\log(z/0.1)}{\log(2000/0.1)}-1.
\end{equation}
Within sky regions, the nearest valid log-depth value is extended for 16 target-resolution pixels.
From 16 to 64 pixels, a smoothstep blend transitions to the far-depth value $+1$, which is used beyond 64 pixels and for all-sky frames.
The resulting depth maps are replicated across three channels and encoded by the frozen Wan VAE, preserving temporal correspondence with all 81 RGB frames.

\subsection{Depth evaluation protocol}

The depth comparison uses a separate set of 16 GTA clips, with all 81 frames evaluated at $832\!\times\!480$.
We compute each metric over all valid pixels in a clip, then average equally across clips.
AbsRel is the mean absolute depth error divided by reference depth, and RMSE is reported in metres.
Predictions are evaluated at their inferred scale.

For temporal evaluation, reference depth, intrinsics, and poses project pixels into the next frame on a four-pixel sampling grid.
We retain visible correspondences with reference depths in $[0.1,1000)$\,m and a depth-consistency tolerance of $\max(0.5\,\mathrm{m},0.03z)$.
The temporal error is the mean absolute difference between predicted log-depth change and reference log-depth change along these motion-compensated correspondences.
Reference geometry supplies the evaluation correspondences and visibility checks.

Table~\ref{tab:supp_depth} compares external focal input and VDA stabilization.
The selected RGB-only combination has the lowest temporal log-depth error among these configurations.
MetricAnything alone without external focal input has lower per-frame AbsRel and RMSE over 0.1--80\,m.
The selected combination reduces temporal error by 43.5\%, with an AbsRel increase of 0.51 percentage points and an RMSE increase of 0.09\,m in this range.

% Camera details and ablations follow camera_trajectory_report.md.
% Depth inference follows df-full-20261102/wan/scope/metric_depth_*.py.
% Depth comparison follows metric-depth-rgb-v1.0/evidence/nok_results.csv,
% protocol.json, and nok_protocol.json. This panel contains 16 GTA clips.

\section{Video Evaluation Details}
\label{sec:supp_video_evaluation}

The comparison uses 1,350 clips selected from the data sources described in the main paper, with 81 frames at 16\,fps.
Our model and SolarWM use four denoising steps at $832\!\times\!480$ and $864\!\times\!480$, respectively.
Pixel-level references follow the corresponding native preprocessing, and minWM retains its native camera-translation normalization.

PSNR, SSIM, and LPIPS with AlexNet features~\cite{zhang2018lpips} exclude the conditioning frame and are averaged equally across clips.
FVD~\cite{unterthiner2018fvd} uses all 81 frames and a shared 400-dimensional pre-softmax I3D representation at $224\!\times\!224$ resolution.
ViPE~\cite{huang2025vipevideoposeengine} recovers trajectories from generated RGB for camera-control evaluation.
We report SCoPE's framewise RotErr in degrees, TransErr, and CamMC~\cite{yin2026scope}, together with pose-matrix cosine similarity and all-pairs relative-pose AUC at $15^\circ$ and $30^\circ$~\cite{wang2025vggt}.
Trajectories are anchored at the first frame and translations are independently normalized by maximum displacement, giving relative-motion scores.
